\documentclass[a4paper,12pt]{article}
\usepackage[T2A]{fontenc}
\usepackage[utf8]{inputenc}
\usepackage[russian]{babel}
\usepackage{amsmath,amssymb,mathtools}
\usepackage{helvet}
\renewcommand{\familydefault}{\sfdefault}
\usepackage{icomma}
\usepackage[a4paper,left=22mm,right=22mm,top=22mm,bottom=23mm]{geometry}
\usepackage{microtype}
\usepackage{tikz}
\usetikzlibrary{arrows.meta,positioning,shapes.geometric,calc}
\usepackage{tabularx,booktabs}
\usepackage{enumitem}
\usepackage{hyperref}
\usepackage{parskip}
\usepackage{fancyhdr}
\usepackage{setspace}
\usepackage{xcolor}
\usepackage{needspace}
\setlength{\headheight}{25pt}
\definecolor{accent}{HTML}{167478}
\definecolor{darkaccent}{HTML}{16474C}
\definecolor{textgray}{HTML}{343B3C}
\definecolor{softline}{HTML}{B8CDCB}
\hypersetup{colorlinks=true,linkcolor=darkaccent,urlcolor=accent,pdfauthor={Артём Александрович Лёвин},pdftitle={Чек-лист: свойства степеней и показательные уравнения}}
\setstretch{1.13}
\setlength{\parindent}{0pt}
\setlength{\parskip}{5pt plus 1pt minus 1pt}
\setlist[enumerate]{itemsep=6pt,topsep=5pt,leftmargin=*,label=\textcolor{accent}{\bfseries\arabic*.}}
\setlist[itemize]{itemsep=2pt,topsep=3pt,leftmargin=*,label=\textcolor{accent}{\textbullet}}
\pagestyle{fancy}
\fancyhf{}
\fancyhead[L]{\small\textcolor{darkaccent}{МАТЕМАТИКА \quad/\quad ЕГЭ, профиль}}
\fancyhead[R]{\small\textcolor{darkaccent}{19.10.26}}
\fancyfoot[C]{\small\textcolor{textgray}{\thepage}}
\renewcommand{\headrulewidth}{0.3pt}
\renewcommand{\headrule}{\hbox to\headwidth{\color{softline}\leaders\hrule height \headrulewidth\hfill}}
\newcommand{\tangent}{\operatorname{tg}}
\newcommand{\key}[1]{\textbf{\textcolor{darkaccent}{#1}}}
\newcommand{\attention}[1]{\medskip\noindent\textcolor{darkaccent}{\bfseries Важно.}\ #1\par\medskip}
\begin{document}

% Титульный лист
\thispagestyle{empty}
\vspace*{34mm}
\begin{center}
{\large\textcolor{accent}{МАТЕМАТИКА \quad\textbullet\quad ЕГЭ ПРОФИЛЬНЫЙ}}\par
\vspace{17mm}
{\fontsize{35}{40}\selectfont\bfseries\textcolor{darkaccent}{Чек-лист}}\par
\vspace{12mm}
{\LARGE\bfseries Свойства степеней\\[5pt] и показательные уравнения}\par
\vspace{7mm}
{\large С повторением тригонометрических преобразований}\par
\vspace{19mm}
{\Large Ярослав Гаврилов}\par
\vspace{5mm}
{\large 19 октября 2026 года}\par
{\normalsize\textcolor{textgray}{По материалам занятия 18 октября 2026 года}}\par
\end{center}
\vfill
\begin{center}
\textcolor{darkaccent}{Лёвин Артём Александрович}\par
\textcolor{textgray}{эксклюзивно для Ярослава Гаврилова}
\end{center}
\begin{tikzpicture}[remember picture,overlay]
\draw[accent,line width=1.15pt,opacity=0.7]
  ([xshift=14mm,yshift=10mm]current page.south west)
  .. controls ([xshift=61mm,yshift=19mm]current page.south west)
           and ([xshift=132mm,yshift=6mm]current page.south west)
  .. ([xshift=-14mm,yshift=12mm]current page.south east);
\end{tikzpicture}
\clearpage

\section*{01 \quad Тригонометрия: приводим углы к одному значению}
\key{Главная идея.} Если в дроби присутствуют тригонометрические функции разных углов, ищем связь через $180^\circ$ и формулы приведения.
\[
\begin{aligned}
\tangent(180^\circ-\alpha)&=-\tangent\alpha,
&\tangent(180^\circ+\alpha)&=\tangent\alpha,\\
\sin(180^\circ-\alpha)&=\sin\alpha,
&\sin(180^\circ+\alpha)&=-\sin\alpha.
\end{aligned}
\]
\textbf{Пример 1.} Упростим выражение:
\[
\frac{5\tangent163^\circ}{\tangent17^\circ}
=\frac{5\tangent(180^\circ-17^\circ)}{\tangent17^\circ}
=\frac{-5\tangent17^\circ}{\tangent17^\circ}=\boxed{-5}.
\]
\textbf{Пример 2.} Сначала замечаем формулу двойного угла $\sin2\alpha=2\sin\alpha\cos\alpha$:
\[
\begin{aligned}
\frac{10\sin85^\circ\cos85^\circ}{\sin190^\circ}
&=\frac{5\sin170^\circ}{\sin190^\circ}
=\frac{5\sin(180^\circ-10^\circ)}{\sin(180^\circ+10^\circ)}\\
&=\frac{5\sin10^\circ}{-\sin10^\circ}=\boxed{-5}.
\end{aligned}
\]
\attention{Сокращать можно только ненулевой множитель. Для указанных углов $\sin10^\circ\ne0$ и $\tangent17^\circ\ne0$.}

\section*{02 \quad Свойства степеней}
Для действительных показателей степеней все приведённые ниже равенства гарантированно верны при \key{$a>0$ и $b>0$}. В конкретных задачах с целыми степенями условия могут быть шире.
\renewcommand{\arraystretch}{1.35}
\begin{center}
\small
\begin{tabularx}{\linewidth}{@{}p{0.45\linewidth}X@{}}
\toprule
\textbf{Равенство} & \textbf{Когда применять} \\
\midrule
$a^m\cdot a^n=a^{m+n}$ & Умножение степеней с одинаковым основанием. \\
$\dfrac{a^m}{a^n}=a^{m-n}$ & Деление степеней с одинаковым основанием. \\
$a^n b^n=(ab)^n$ & Одинаковый показатель; произведение оснований. \\
$\dfrac{a^n}{b^n}=\left(\dfrac ab\right)^n$ & Одинаковый показатель; частное оснований. \\
$(a^m)^n=a^{mn}$ & Степень возводится в другую степень. \\
$\dfrac1{a^n}=a^{-n}$ & Перевод обратной величины в отрицательную степень. \\
$\left(\dfrac ab\right)^n=\left(\dfrac ba\right)^{-n}$ & Обращение дроби меняет знак показателя. \\
$\sqrt[k]{a^m}=a^{m/k}$, $k\ge2$ & Корень записывается дробным показателем. \\
\bottomrule
\end{tabularx}
\end{center}
\attention{Формулы работают в обе стороны, но $a^m+a^n\ne a^{m+n}$ в общем случае. Складывать показатели разрешено именно при \emph{умножении} степеней с одинаковым основанием.}

\clearpage
\section*{03 \quad Общая стратегия: привести к степенному виду}
\key{Порядок действий:} составные числа $\rightarrow$ корни $\rightarrow$ дроби $\rightarrow$ свойства степеней.
\begin{enumerate}
\item \textbf{Разложите составные числа} на простые множители: $8=2^3$, $81=3^4$, $625=5^4$, $3125=5^5$.
\item \textbf{Замените корни} дробными показателями: $\sqrt8=\sqrt{2^3}=2^{3/2}$; $\sqrt[5]{81}=3^{4/5}$.
\item \textbf{Представьте подходящие дроби} через отрицательные показатели: $\dfrac1{\sqrt8}=2^{-3/2}$; $0{,}2=5^{-1}$; $0{,}5=2^{-1}$.
\item \textbf{Приведите основания к общему виду}, затем сложите, вычтите или перемножьте показатели по соответствующему правилу.
\item \textbf{Проверьте скобки, знаки и область допустимых значений.}
\end{enumerate}
\textbf{Пример.} Последовательно приводим к одному основанию:
\[
\frac{1}{\sqrt8}=
\frac1{\sqrt{2^3}}=\frac1{2^{3/2}}=\boxed{2^{-3/2}}.
\]
\textbf{Важное замечание об обозначениях.} У любого основания без указанного показателя есть подразумеваемая первая степень: $2=2^1$. Обычный квадратный корень имеет индекс $2$: $\sqrt2=\sqrt[2]{2}=2^{1/2}$.

\textbf{Выбор формулы.} Одинаковые основания --- складываем или вычитаем показатели; одинаковые показатели --- объединяем основания; степень в степени --- перемножаем показатели.

\textbf{Проверка смысла:} $(a+b)^n\ne a^n+b^n$ в общем случае; $\sqrt{x^2}=|x|$ для действительного $x$.

\section*{04 \quad Вычисления по законам степеней}
\textbf{Пример 1.} Основания сначала делаем одинаковыми:
\[
25^{0{,}18}\cdot5^{0{,}64}
=(5^2)^{0{,}18}\cdot5^{0{,}64}
=5^{0{,}36+0{,}64}=\boxed5.
\]
\textbf{Пример 2.} Корень переводим в показатель степени:
\[
\begin{aligned}
\left(\frac{2^{1/3}\cdot2^{1/4}}{\sqrt[12]{2}}\right)^2
&=\left(2^{1/3+1/4-1/12}\right)^2
=\left(2^{1/2}\right)^2=\boxed2.
\end{aligned}
\]
\textbf{Пример 3.} Десятичная дробь и правило скобок:
\[
\begin{aligned}
\frac{(0{,}5)^{\sqrt{10}-1}}{2^{-\sqrt{10}}}
&=2^{-(\sqrt{10}-1)-(-\sqrt{10})}\\
&=2^{-\sqrt{10}+1+\sqrt{10}}=\boxed2.
\end{aligned}
\]
\attention{Правило «ВУПС»: \textbf{в}ычитаемое, \textbf{у}множаемое и \textbf{п}одставляемое --- \textbf{с}кобки. Особенно осторожно записывайте вычитание отрицательных показателей.}

\clearpage
\section*{05 \quad Алгоритм преобразования степенных выражений}
\vspace{7mm}
\begin{center}
\begin{tikzpicture}[
 every node/.style={font=\small,align=center},
 startstop/.style={ellipse,draw=darkaccent,very thick,minimum width=43mm,minimum height=10mm,fill=white},
 process/.style={rectangle,rounded corners=2pt,draw=softline,line width=0.8pt,minimum width=119mm,text width=108mm,minimum height=14mm,fill=white},
 arr/.style={-{Stealth[length=2.6mm]},thick,draw=accent},
 node distance=11mm
]
\node[startstop] (st) {Начало};
\node[process,below=of st] (a) {Если есть составные числа --- разложить их на простые множители};
\node[process,below=of a] (b) {Если есть корни --- заменить их дробными степенями};
\node[process,below=of b] (c) {Если есть дроби --- использовать отрицательные степени, где это удобно};
\node[process,below=of c] (d) {Привести подходящие выражения к одинаковым основаниям};
\node[process,below=of d] (e) {Применить свойства степеней и упростить показатели};
\node[process,below=of e] (f) {Проверить скобки, знаки и допустимость преобразований};
\node[startstop,below=of f] (end) {Ответ};
\draw[arr] (st)--(a);
\draw[arr] (a)--(b);
\draw[arr] (b)--(c);
\draw[arr] (c)--(d);
\draw[arr] (d)--(e);
\draw[arr] (e)--(f);
\draw[arr] (f)--(end);
\end{tikzpicture}
\end{center}
\vfill
\begin{center}
\small\textcolor{textgray}{Каждый шаг выполняется только при наличии соответствующего объекта в выражении.}
\end{center}
\clearpage

\section*{06 \quad Показательные уравнения}
\key{Основной принцип.} Если $a>0$, $a\ne1$, то функция $a^t$ строго монотонна, поэтому
\[
\boxed{\quad a^{f(x)}=a^{g(x)}\iff f(x)=g(x).\quad}
\]
Перед приравниванием показателей в каждой части должна стоять \emph{одна степень с одинаковым основанием}, без посторонних множителей и слагаемых.
\begin{enumerate}
\item $2^x=8=2^3\ \Rightarrow\ \boxed{x=3}$.
\item $2\cdot2^x=8\ \Rightarrow\ 2^{x+1}=2^3\ \Rightarrow\ \boxed{x=2}$.
\item $40\cdot2^{-x/3}=5\ \Rightarrow\ 2^{-x/3}=\dfrac18=2^{-3}\ \Rightarrow\ \boxed{x=9}$.
\item $16\cdot2^{-x/5}=1\ \Rightarrow\ 2^{-x/5}=2^{-4}\ \Rightarrow\ \boxed{x=20}$.
\end{enumerate}
\subsection*{Как убрать дробный или отрицательный показатель}
Если $x\ge0$ и $x^{5/3}=32$, возводим \emph{обе части} в степень $3/5$:
\[
x=(32)^{3/5}=(2^5)^{3/5}=\boxed8.
\]
Аналогично, при $x\ge0$ имеем $x^{3/2}=8\Rightarrow x=8^{2/3}=\boxed4$.

\attention{Для уравнения $x^{-2}=4$ на всей области $x\in\mathbb R\setminus\{0\}$ получаются \textbf{два корня}: $x=\pm\frac12$. Если дано дополнительное условие $x>0$, остаётся только $x=\frac12$. При чётных степенях нельзя автоматически терять отрицательный корень.}

\section*{07 \quad Прикладная задача: степенная зависимость}
\textbf{Условие.} Величины $P$ и $v>0$ связаны соотношением
$Pv^{5/3}=C$. Найдите $v$, если $P=6{,}25\cdot10^5$, $C=200$.
\[
\begin{aligned}
Pv^{5/3}&=C, & v^{5/3}&=\frac CP=\frac{200}{625000}=\frac1{3125}=5^{-5},\\
v&=\left(5^{-5}\right)^{3/5}=5^{-3}=\boxed{\frac1{125}=0{,}008}.
\end{aligned}
\]
\textbf{Порядок решения:} выписать данные $\rightarrow$ подставить в формулу $\rightarrow$ изолировать степень неизвестной $\rightarrow$ привести числа к общему основанию $\rightarrow$ возвести обе части в обратную степень $\rightarrow$ проверить результат и единицы измерения.

\clearpage
\section*{08 \quad Тренировочный блок \textnormal{(Путь А, 10 заданий)}}
\textcolor{textgray}{Решите самостоятельно. Сначала определите применимое правило; ответы размещены на следующей странице.}
\begin{enumerate}
\item Вычислите $\displaystyle\frac{7\tangent152^\circ}{\tangent28^\circ}$.
\item Вычислите $\displaystyle\frac{8\sin70^\circ\cos70^\circ}{\sin220^\circ}$.
\item Представьте в виде степени с основанием $2$ число $\displaystyle\frac1{\sqrt[5]{32}}$.
\item Упростите до одной степени с основанием $3$: $\displaystyle 81^{1/4}\cdot3^{1/2}$.
\item Вычислите $\displaystyle16^{0{,}375}\cdot2^{0{,}5}$.
\item Вычислите $\displaystyle\left(\frac{3^{1/2}\cdot3^{1/3}}{\sqrt[6]{3}}\right)^3$.
\item Вычислите $\displaystyle\frac{(0{,}25)^{\sqrt5-2}}{2^{-2\sqrt5}}$.
\item Решите уравнение $\displaystyle 27\cdot3^{(x-2)/3}=9$.
\item Решите уравнение $\displaystyle x^{3/2}=64$ (действительные числа).
\item \textbf{Прикладная задача.} Величины $P$ и $v>0$ связаны формулой $Pv^{3/2}=C$. Найдите $v$, если $P=2{,}5\cdot10^3$, $C=20$. Ответ запишите десятичной дробью.
\end{enumerate}
\vfill
\textcolor{darkaccent}{\textbf{Самопроверка перед сдачей:}}\quad одинаковые основания; корректные показатели; скобки при вычитании; ограничения для степеней; проверка подстановкой.
\clearpage

\section*{09 \quad Ответы к тренировочному блоку}
\renewcommand{\arraystretch}{1.55}
\begin{center}
\begin{tabularx}{\linewidth}{@{}c X@{}}
\toprule
\textbf{№} & \textbf{Ответ / ключевое преобразование}\\
\midrule
1 & $-7$; $152^\circ=180^\circ-28^\circ$. \\
2 & $-4$; $8\sin70^\circ\cos70^\circ=4\sin140^\circ$. \\
3 & $2^{-1}$. \\
4 & $3^{3/2}$; $81=3^4$. \\
5 & $4$; показатель при основании $2$ равен $4\cdot0{,}375+0{,}5=2$. \\
6 & $9$; $3\left(\frac12+\frac13-\frac16\right)=2$. \\
7 & $16$; $(0{,}25)^{\sqrt5-2}=2^{-2(\sqrt5-2)}$. \\
8 & $x=-1$. \\
9 & $x=16$. \\
10 & $v=0{,}04$. \\
\bottomrule
\end{tabularx}
\end{center}
\medskip
\textbf{Проверка двух наиболее сложных упражнений.}
\[
\begin{aligned}
\text{№ 8: }&\quad 3^{3+(x-2)/3}=3^2
\ \Rightarrow\ 3+\frac{x-2}{3}=2
\ \Rightarrow\ x=-1;\\[6pt]
\text{№ 10: }&\quad v^{3/2}=\frac{20}{2500}=\frac1{125}=5^{-3},\\
&\quad v=(5^{-3})^{2/3}=5^{-2}=\frac1{25}=0{,}04.
\end{aligned}
\]
\vfill
\begin{center}
\textcolor{darkaccent}{\small Лёвин Артём Александрович \quad\textbullet\quad эксклюзивно для Ярослава Гаврилова}
\end{center}
\end{document}
